% =====================================================
% Plantilla adaptada de FloripaSat-2 / GOLDS-UFSC §4.4
% Para usar en: 00_documento_principal/chapters/02_sistema.tex (§2.2.3)
% Fuente de datos: 03_recursos/presupuestos/Requerimiento de energía - INTISAT.xlsx
% =====================================================

\subsection{Presupuesto de Potencia (Power Budget)}

El presupuesto de potencia se construye en tres pasos (ref. \cite{Wertz2011}):
\begin{enumerate}
    \item Preparar el \textbf{presupuesto de potencia operacional} (consumo por subsistema y modo).
    \item Dimensionar la \textbf{batería} (DoD, ciclos, profundidad de descarga).
    \item Estimar la \textbf{degradación} de potencia sobre la vida útil de la misión.
\end{enumerate}

% -----------------------------------------------------
\subsubsection*{Entrada de potencia (paneles solares)}

\begin{table}[h]
\centering
\caption{Potencia generada por los paneles solares (simulación orbital).}
\label{tab:power_input}
\begin{tabular}{lr}
\toprule
\textbf{Parámetro} & \textbf{Valor} \\
\midrule
Potencia pico (sin eclipse)         & TBD~mW \\
Potencia promedio (con eclipse)     & TBD~mW \\
Duración media iluminación (sunlight) & TBD~min \\
Duración media eclipse              & TBD~min \\
\bottomrule
\end{tabular}
\end{table}

% -----------------------------------------------------
\subsubsection*{Consumo operacional por subsistema}

Tomando como fuente el archivo
\texttt{03\_recursos/presupuestos/Requerimiento de energía - INTISAT.xlsx}:

\begin{table}[h]
\centering
\caption{Consumo de potencia por módulo y \textit{duty cycle}.}
\label{tab:power_consumption}
\begin{tabular}{lrr}
\toprule
\textbf{Módulo} & \textbf{Duty Cycle [\%]} & \textbf{Potencia [mW]} \\
\midrule
OBC (rail 3.3 V)               & 100  & 4060 \\
OBC (rail 5 V, CANBUS)         & 100  & 1100 \\
TTC UHF (RX)                   & 95   & TBD \\
TTC UHF (TX, beacon)           & 5    & TBD \\
EPS (idle)                     & 100  & TBD \\
Batería (heater, eclipse)      & 10   & TBD \\
Antena (despliegue, una vez)   & 0    & TBD \\
Payload Microscopía CMOS       & TBD  & TBD \\
Payload IA / SAI (planificado) & 0    & --- \\
\midrule
\textbf{Total satélite}        & ---  & \textbf{TBD} \\
\bottomrule
\end{tabular}
\end{table}

% -----------------------------------------------------
\subsubsection*{Dimensionamiento de batería}

Pasos (ref. \cite{Wertz2011}, sec.~10.3):
\begin{enumerate}
    \item Estimar nivel de potencia que debe suministrar la batería.
    \item Calcular duración de descarga, carga y número de ciclos.
    \item Seleccionar profundidad de descarga (DoD).
    \item Seleccionar tasa de carga.
    \item Calcular potencia de recarga.
\end{enumerate}

\paragraph{Profundidad de descarga (DoD).}
\begin{equation}
\mathrm{DoD} = \frac{D_i \cdot D_t}{C} \cdot 100\,\%
\label{eq:dod}
\end{equation}
donde $D_i$ es la corriente de descarga, $D_t$ el periodo de descarga y $C$ la capacidad de la batería.

\begin{table}[h]
\centering
\caption{Resultados del dimensionamiento de batería.}
\label{tab:battery_sizing}
\begin{tabular}{lrr}
\toprule
\textbf{Parámetro} & \textbf{Valor} & \textbf{Unidad} \\
\midrule
Capacidad de batería                 & TBD & mWh \\
Consumo en iluminación               & TBD & mW \\
Consumo en eclipse                   & TBD & mW \\
Duración iluminación                 & TBD & h \\
Duración eclipse                     & TBD & h \\
DoD calculado                        & TBD & \% \\
Margen energético                    & TBD & mWh \\
\bottomrule
\end{tabular}
\end{table}

\paragraph{Criterio de aceptación.}
Para una órbita LEO con 5 años de vida útil (típico CubeSat), el DoD recomendado es
$\le 25\,\%$. FloripaSat-2 cerró con DoD$=8{,}5\,\%$ y margen energético positivo (+643~mWh por órbita).
